\scanpage{10}
\textcircled{2}\quad Supposons que $U$ soit \deleted{muni d'une} \inserted{une \textcircled{$k$}-algèbre}
\margin{$k$}
augmentée ($k$ anneau comm.) avec un
hom diagonal
\[
\Delta:U\unsure{\text{—}|\longrightarrow} U\mathbin{\widehat\otimes}_k U
\]
qui soit un hom de $k$-algèbres augmentées.
$U$ est appelé une $k$-hyperalgèbre \inserted{topologique} (au sens
large). On désigne par $\deleted{\unsure{G(U)|\pi(U)}}\deleted{\unsure{gr(U)|\mathfrak g(U)}}\unsure{\gamma|\mathscr G}(U)$ l'ens
des $g\in 1+I(U)$ \deleted{tels que} tels que
\[
\begin{aligned}
\Delta(g)&=g\otimes g\\
\delta\xi&=\xi\otimes\xi\qquad\text{si}\quad g=1+\xi\quad(\xi\in I(U))
\end{aligned}
\]
\margin{on suppose que $I(U)^n\to0$ pour $n\to\infty$ \deleted{\unsure{Donc}}}
[\deleted{\unsure{éléments}}\inserted{i.e.} ils correspondent par dualité aux
hom. de la hyperalgèbre \deleted{\unsure{triviale}}
$k[T]$, avec $\delta T=T\otimes T$ ; on peut aussi
prendre $k[[T]]$, avec $\delta T=T\otimes T$, et
se borner aux hom. \unsure{continus|continues} $\ldots$ \unsure{]|\}} ;
cette \inserted{hyper}algèbre est le dual top. de
l'algèbre du groupe $\mathbb Z$ $\ldots$].
\deleted{Ils} Ces $g$ forment un groupe, car
\[
\begin{aligned}
\Delta(gg')&=\Delta(g)\Delta(g')=(g\otimes g)(g'\otimes g')=gg'\otimes g\unsure{g'}\\
\Delta(g^{-1})&=\deleted{\unsure{g\otimes}}(\Delta(g))^{-1}=(g\otimes g)^{-1}=g^{-1}\otimes g^{-1}.
\end{aligned}
\]
Fonctorialité de $\gamma(U)$, et $\pi(U)$ (primitifs)
\deleted{\unsure{pour}} \deleted{\unsure{les|en}} $(k,U)\longrightarrow(k',U')$.
% archive sheet number: 9

